\documentclass[tikz, border=5mm]{standalone}
\usepackage{geometricfigures}
\usepackage{amsmath}
\usepackage{amsfonts}
\usepackage{amssymb}
\usetikzlibrary{arrows, positioning, shapes.geometric, calc}

\begin{document}
\begin{tikzpicture}[
    scale=0.9, transform shape,
    torus/.style={ellipse, draw, fill=mgreen!20!bg, minimum width=4cm, minimum height=2.8cm} % Torus style
]
    % --- Lie Group G (Drawn as a Torus Node) ---
    \node[torus, label={[name=G_label, yshift=-1.2cm, xshift=-1.7cm]$G$}] (G_node) at (0,3) {}; % Renamed to G_node to avoid conflict
    \draw[fill=bg] (G_node.center) ellipse (0.8cm and 0.4cm);

    % Define a point on the "torus" to be our identity 'e'
    \coordinate (e_on_G) at ($(G_node.center) + (0.7, -0.5)$);
    \fill[mred] (e_on_G) circle (2pt) node[xshift=.5em] {$\mathbb{I}$};

    % A point 'g' near 'e' on the Lie group
    \coordinate (g_on_G) at ($(G_node.center) + (0.5, 0.4)$);
    \fill[orange] (g_on_G) circle (1.5pt);
    \node[orange, above right=1pt of g_on_G] {$g$};

    % Tangent Space (Lie Algebra)
    % Position it below the identity, roughly aligned
    \coordinate (below_e) at ($(e_on_G) + (0.0, -2.0)$);
    \draw[thick, muted, fill=muted!20!bg, opacity=0.8] ($(below_e) + (-1.0, 0.0)$) -- ($(below_e) + (2.5, -0.0)$) -- ($(below_e) + (1.5, -2.5)$) -- ($(below_e) + (-2.0, -2.5)$) -- cycle;
    \node[muted] at ($(below_e) + (0.5, -2.2)$) {$\mathfrak{g} \simeq T_\mathbb{I} G$};

    % Example element in Lie Algebra
    \coordinate (arrow_tip) at ($(below_e) + (0.5, -1.5)$);
    \draw[->, thick, fg] (below_e) -- (arrow_tip);
    \node[fg, below right] at ($(e_on_G) + (0, -1)$) {$X \in \mathfrak{g}$};

    % Exponential map
    \draw[->, very thick, mred, rounded corners] ($(arrow_tip) + (0.03, 0.0)$) .. controls ($(below_e) + (1.5, 0.5)$) and ($(below_e) + (1.5, 0.5)$) .. (g_on_G) node[midway, xshift=1em] {$\exp$};

    % Connecting Lines
    \draw[dashed, muted] (e_on_G.south) -- (below_e); % Connect identity to origin of tangent space
\end{tikzpicture}
\end{document}
